% AUTO-GENERATED by scripts/build_topics.py — do not edit.
% Edit topics/bodies/topic_6.tex or topics/preamble.tex instead.
%% Placeholder. Replace with the written chapter.
\textit{Източници:} снимки \texttt{6та тема}

\begin{supp}[Бележка към картографирането]
6.pdf is STALE — it holds 'Модели на изчисленията: машина на Тюринг, RAM', a question from an earlier конспект. The photographs (двоична пирамида, HEAPSORT) are correct. Do not use 6.pdf for this chapter.
\end{supp}

\subsection*{Анотация от конспекта}
Сортиране чрез сравнения във време O(n lg n). Двоична пирамида (binary heap): определение и свойства. Наивно построяване на двоична пирамида: доказателство за коректност и изследване на сложността по време. Бързо построяване на двоична пирамида, използващо функция Heapify: доказателство за коректност и изследване на сложността по време. Сортиращ алгоритъм HEAPSORT: псевдокод, доказателство за коректност и изследване на сложността по време. Сортиращ алгоритъм MERGESORT като пример за алгоритъм по схемата Разделяй и Владей. Псевдокод и доказателство за коректност на MERGESORT. Изследване на сложността по време на MERGESORT. Приложение на MERGESORT за намиране на броя на инверсиите в масив от числа, с доказателство за коректност. Забележка: На изпита ще се падне или HEAPSORT, или MERGESORT. Литература: [31].
